\section{Related Work}
\label{sec:related}

\paragraph{Scaling laws for compression, in loss space.}
Established compression laws are laws of pretraining loss:
sparsity acts as an effective-parameter multiplier
\citep{frantar2023scaling}, precision as a saturating effective-parameter
count with a post-training degradation term \citep{kumar2024scaling,
dettmers2023case}, and sparse/quantized formats unify under a single
capacity scalar \citep{panferov2025unified}; distillation obeys a broken
power law in teacher loss \citep{busbridge2025distillation}; pruned models
admit post-training recovery laws \citep{zhang2024scaling}, and pruning
plus distillation-based recovery is highly effective in practice
\citep{muralidharan2024compact,xia2023sheared}. None of these validate
their forms on downstream capabilities. Empirical claims of
``compression laws'' with downstream metrics exist
\citep{sengupta2025compression} but report no held-out predictive
validation and no capability stratification; task-stratified PTQ laws
\citep{zhou2025task} are the closest antecedent, limited to quantization
without a mechanism or recovery axis. Our framework is complementary: it
derives per-capability forms from measured geometry and validates them by
extrapolation.

\paragraph{Predictability of downstream capabilities.}
Downstream accuracy is predictable essentially only through continuous
log-likelihood measurements: emergent jumps are largely metric artifacts
\citep{schaeffer2023emergent}, accuracy-space prediction degrades through a
chain of information-destroying transforms \citep{schaeffer2024predicting},
and successful pipelines fit loss-space power laws composed with
sigmoid/exponential links \citep{gadre2024language,bhagia2024establishing},
aggregate many tasks \citep{owen2024how}, or use multi-model observational
latents \citep{ruan2024observational}. Even then, linear loss-to-downstream
relationships hold in only a minority of settings \citep{lourie2025scaling}.
Our measurement layer adopts these lessons wholesale
(\S\ref{sec:prelim}).

\paragraph{Capability-selective effects of compression.}
Pruning disproportionately harms knowledge-intensive tasks
\citep{jaiswal2023compressing}, low-bit quantization disproportionately
harms reasoning \citep{liu2025quantization}, small-magnitude weights encode
hard-task knowledge whose removal is not recovered by further training
\citep{yin2023junk}, and selective low-rank truncation can \emph{improve}
specific behaviors \citep{sharma2023laser}. These observations motivate a
capability-vector law and anticipate our signed first-order term; what has
been missing is a quantitative, predictive account, which capability
geometry supplies.

\paragraph{Curvature, gradients, and capability structure.}
Second-order weight importance underlies classical and modern pruning
\citep{lecun1989optimal,hassibi1992second,frantar2023sparsegpt,
sun2023simple}, and Fisher-aligned subspaces improve compression choices
\citep{fasc2026}; per-task Fisher guidance exists in vision transformers
\citep{fisher3d2026}. Gradient structure also localizes skills: gradients
occupy low-dimensional subspaces \citep{gurari2018gradient,zhao2024galore},
gradient clustering discovers skill quanta \citep{michaud2023quantization},
and task vectors localize capabilities in weight space
\citep{ilharco2023editing}. To our knowledge, no prior work connects
capability-conditional gradient/Fisher structure of the \emph{source} model
to cross-method compression damage and recovery; this is the gap the
present paper occupies.

\paragraph{Origins of scaling-law forms.}
We follow the tradition that credible scaling-law forms carry mechanism:
error decompositions \citep{hoffmann2022training}, solvable models and
data-spectrum arguments \citep{sharma2020neural,bahri2021explaining,
maloney2022solvable}, discrete-skill models \citep{michaud2023quantization},
and axiomatic limit constraints \citep{busbridge2025distillation}; surveys
document the fragility of purely fitted forms \citep{li2025misfitting,
besiroglu2024chinchilla}. Our two-regime law with geometry-constrained
coefficients is designed to meet this bar.
